% ---------------------------------------------------------------------
% Section 8
% ---------------------------------------------------------------------

\begin{cnread}
两个带电粒子「带着动量碰撞」在 \AEG{} 里对应什么？本节给出一个纯代数的
区分：$e^-e^-$ 散射（M\o ller）只有 $t$-道与 $u$-道，两者都留在仿射扇区
$B_\infty$ 内；而 $e^+e^-$ 散射（Bhabha）多出一条 $s$-道，它在 \AEG{} 中
恰好是\textbf{跨越 Bruhat 大胞腔}：$\PGL_2=B_\infty\sqcup B_\infty JB_\infty$
（Paper~0 公式 47），必经一次反演 $J(z)=-1/z$。总加法荷归零（$+e-e=0$）
正是「分母归零」的代数来源，而射影延拓让人从 $z=0$ 穿到 $\infty$。
这里要提醒：这是\textbf{结构对应}，不是振幅的推导。
\end{cnread}

\subsection{Two electrons: both channels inside the affine sector}

Let both particles be electrons.  Then the total additive charge of any history
containing both is $-2e\ne0$, so no closed charge loop and hence no annihilation
channel exists: by Proposition~\ref{prop:dipole-loop} the closed-loop mechanism
requires opposite charges.

Two histories remain, and both are affine.

\begin{enumerate}[label=(\roman*)]
  \item \emph{Threading order.}  The word exchanges one quantum while the two
  world lines keep their order in the scale coordinate; this is the
  ``$t$-channel'' history of the exchange, and it is a word in
  $B_\infty$ whose total charge is $(-e,\,-e)$.
  \item \emph{Crossed order.}  The two final-state world lines are exchanged, so
  the appropriate comparison is between two histories $\gamma$ and $\delta$
  whose scale steps are interleaved; this is the ``$u$-channel'' of identical
  particles, and by Theorem~8.9 of Paper~0 the relative defect is the weighted
  area between the two charge chains.
\end{enumerate}

In both cases the surviving velocity dependence is the one computed in
Proposition~\ref{prop:ampere}: the two histories differ in their direction angle
$\theta_1-\theta_2$, and the pairing gives the coefficient
$1-\mathbf v_1\!\cdot\!\mathbf v_2/c^{2}$ relative to the Coulomb term.  For
collinear motion this is the correction computed in
Corollary~\ref{cor:parallel-currents}.

\begin{remark}[On spin]
\label{rem:spin}
A spin-$\tfrac12$ charge would need the double cover of the projective group, i.e.
the passage from $\PSL_2(\C)$ to $\SL_2(\C)$, whose two-dimensional
representation is generated by the Pauli matrices.  Nothing in Paper~0 or
Paper~IV supplies that cover, and no spin-dependent statement is proved in this
note.  The remark is recorded because the physical narrative contains it.
\end{remark}

\subsection{Electron and positron: the second Bruhat cell}

Now let the particles be an electron and a positron.  Two of the three channels
are affine, and one is not.

\begin{enumerate}[label=(\roman*)]
  \item \emph{Exchange without annihilation.}  The ``$t$-channel'' photon
  exchange is the configuration of Section~\ref{sec:coulomb}: a bound-state
  history in $B_\infty$ with potential $V=-\alpha_e/r$ and defect
  $e(1-e^{-M_*})$ by Proposition~\ref{prop:dipole-loop}.
  \item \emph{Annihilation.}  At high relative momentum the pair acquires a
  channel with no like-charge analogue.  Its algebraic signature is that the
  total additive charge vanishes, $A_{\mathrm{tot}}=(+e)+(-e)=0$, exactly as in
  the closed loop of Proposition~\ref{prop:dipole-loop}.
  \item \emph{Regeneration.}  The word then continues from the vanishing
  denominator, and the projective continuation of Paper~0, Section~5.1 carries
  it through infinity and back to a finite pair.
\end{enumerate}

The third item is where the projective cell is entered, and \AEG{} has an exact
statement about it.

\begin{theorem}[Bruhat decomposition, {\cite[Eq.~(47) and Thm.~6.2]{AEG0}}]
\label{thm:bruhat}
For any field $K$,
\begin{equation}\label{eq:bruhat}
  \PGL_2(K)=B_\infty\sqcup B_\infty J B_\infty,
\end{equation}
and every $F\in\PGL_2(K)$ with nonzero lower-left entry admits the exact
factorization
\begin{equation}\label{eq:factorization}
  F=T_{A/C}\circ D_{(AD-BC)/C^{2}}\circ J\circ T_{D/C},
\end{equation}
i.e. translate, scale, invert, translate.  The first cell consists of operators
whose pole is infinity; the second consists of operators with a finite pole in
the chosen affine chart.
\end{theorem}

\begin{proof}
Paper~0, Theorem~6.2 and its proof; the two cases $C=0$ and $C\ne0$ are
$F=T_{B/D}\circ D_{A/D}$ and \eqref{eq:factorization} respectively.
\end{proof}

The reading of the theorem in the scattering picture is now a translation of its
two displayed statements: a history that stays in $B_\infty$ has a pole at
infinity, and a history that produces a finite pole must contain the inverting
element exactly once.

\begin{center}\small
\renewcommand{\arraystretch}{1.25}
\begin{tabular}{@{}p{0.2\linewidth}p{0.24\linewidth}p{0.24\linewidth}p{0.24\linewidth}@{}}
\toprule
process & charges & \AEG{} sector & mechanism \\
\midrule
$e^-e^-\to e^-e^-$ & $(-e,-e)$, $A_{\mathrm{tot}}=-2e\ne0$
& $B_\infty$ only & threading and crossed (identical-particle) orders; no closed
charge loop \\
$e^+e^-\to e^+e^-$ & $(-e,+e)$, $A_{\mathrm{tot}}=0$
& $B_\infty$ and $B_\infty JB_\infty$ & exchange, plus a channel that must cross
$J$ because the total additive charge vanishes \\
$(e^-e^-)$ emission & $(-e,-e)$ & $B_\infty$, pole jumps
& radiation as in Section~\ref{sec:radiation}: the pole may move inside the
affine cell \\
\bottomrule
\end{tabular}
\end{center}

\subsection{What the correspondence does and does not assert}

\begin{warning}[Structural, not dynamical]
\label{warn:scattering}
The table above is a correspondence between channel types and Bruhat cells.  It
is not a derivation of a scattering amplitude, a cross section, or a Feynman rule
for the annihilation channel.  In particular:
\begin{enumerate}[label=(\alph*)]
  \item The vanishing of $A_{\mathrm{tot}}$ is an exact statement about charges
  (Theorem~8.6); the statement that a physical annihilation requires it is a
  dictionary statement.
  \item The projective continuation through an inadmissible division is exact in
  \AEG{} (Paper~0, Section~5.1) but says nothing about which intermediate states
  are kinematically allowed.  Ordinary division by zero remains inadmissible;
  what the continuation supplies is a value at the pole, not a physical process.
  \item The separation ``$t$-channel inside $B_\infty$, $s$-channel across
  $B_\infty JB_\infty$'' is a two-cell statement.  A quantitative Bhabha cross
  section would require a measure on histories, i.e. the path-integral structure
  of Remark~\ref{rem:prequantum}, which is not constructed here.
\end{enumerate}
\end{warning}

What survives is the qualitative separation, and it is not empty: it explains why
the like-charge and unlike-charge cases must differ in \emph{kind} rather than
only in sign, and it locates the difference in an exact algebraic invariant (the
Bruhat cell of the composite) rather than in the sign of a coupling.
